\subsection{The Dunford--Pettis theorem}

Let \(\mu\) be a positive measure on T. A vectorial measure \(\mathbf{m}\) on T, with values in F, is said to be scalarly of base \(\mu\) (or to have base \(\mu\) scalarly) if, for every \(\mathbf{z}' \in \mathrm{F}'\) , the scalar measure \(\mathbf{z}' \circ \mathbf{m}\) has base \(\mu\) . If a vectorial measure m with values in F has base \(\mu\) , then it has base \(\mu\) scalarly: for, if \(m = f \cdot \mu\) then \(z' \circ m = \langle z', f \rangle \cdot \mu\) for every \(z' \in F'\) . But there exist vectorial measures that are scalarly of base \(\mu\) without having base \(\mu\) (Exer. 17), and, on the other hand, there exist vectorial measures that are not scalarly of base \(\nu\) for any positive measure \(\nu\) ; note however that every majorizable measure m with values in a normed space is scalarly of base \(|m|\) , by virtue of the Lebesgue--Nikodym theorem.

Example. --- Let us take for m the identity mapping of \(\mathcal{K}(\mathrm{T})\) . To say that m is scalarly of base \(\mu\) means that every real measure on T has base \(\mu\) . In particular, the point measure \(\varepsilon_{t}\) ( \(t \in T\) ) must have base \(\mu\) , which requires that \(\mu(\{t\}) > 0\) for every \(t \in T\) , and implies in particular that every compact subset of T is countable.

In this No.~we are going to prove a result that generalizes one of the consequences of the Lebesgue--Nikodym theorem, namely, that the dual of \(\mathrm{L}^{1}(\mu)\) is \(\mathrm{L}^{\infty}(\mu)\) (Ch. V, §5, No.~8, Th. 4), and that gives a sufficient condition for a vectorial measure that is scalarly of base \(\mu\) to have base \(\mu\) .

Let \(\pi\) be the canonical mapping of \(\mathcal{L}^{\infty}(\mu)\) onto \(\mathrm{L}^{\infty}(\mu)\) . We shall say that a linear subspace \(\mathbf{G}\) of \(\mathrm{L}^{\infty}\) has the lifting property if there exists a linear mapping \(\rho\) of \(\mathbf{G}\) into \(\mathcal{L}^{\infty}(\mu)\) (called a lifting of \(\mathbf{G}\) ) such that \(\pi \circ \rho\) is the identity on \(\mathbf{G}\) and \(|\rho(f)(t)| \leqslant \mathrm{N}_{\infty}(f)\) for all \(t \in \mathrm{T}\) and \(f \in \mathrm{G}\) .

One proves that if \(\mu\) is Lebesgue measure on \(\mathbf{R}^n\) , the entire space \(\mathrm{L}^{\infty}(\mathbf{R}^{n},\mu)\) has the lifting property (Exer. 18).

Lemma 2. --- Every separable subspace G of the Banach space \(\mathrm{L}^{\infty}(\mathrm{T},\mu)\) has the lifting property.

By the hypothesis, there exists a countable dense subset H of G that is a linear subspace with respect to the field Q of rational numbers; let \((h_{n})\) be a (countable) basis of H over Q. For every integer n, let \(h_{n}^{\prime}\) be an element of \(L^{\infty}\) such that \(\pi(h_{n}^{\prime}) = h_{n}\) , and let \(\rho^{\prime}\) be the Q-linear mapping of H into \(L^{\infty}\) defined by \(\rho^{\prime}(h_{n}) = h_{n}^{\prime}\) ; it is clear that \(\pi \circ \rho^{\prime}\) is the identity on H. Moreover, for every \(h \in H\) one has \(|\rho^{\prime}(h)(t)| \leqslant N_{\infty}(h)\) except at the points t of a locally negligible set A(h). Let A be the union of the A(h) for \(h \in H\) , which is also locally negligible. For every \(h \in H\) , denote by \(\rho(h)\) the function \(h^{\prime\prime} \in L^{\infty}\) such that \(h^{\prime\prime}(t) = \rho^{\prime}(h)(t)\) if \(t \notin A\) , and \(h^{\prime\prime}(t) = 0\) if \(t \in A\) . It is clear that \(\rho\) is a Q-linear mapping of H into the subspace B of bounded functions in \(L^{\infty}\) , such that \(\pi \circ \rho\) is the identity on H and such that \(|\rho(h)(t)| \leqslant N_{\infty}(h)\) for all \(h \in H\) and \(t \in T\) . Since B is a Banach space for the norm \(\|f\| = \sup_{t \in T} |f(t)|\) (Ch. IV, §5, No.~4, Th. 2),

\(\rho\) may be extended to a continuous R-linear mapping, again denoted \(\rho\) , of G into B, which is obviously a lifting of G.

DEFINITION 5. --- Let F be a Hausdorff locally convex space, \(F_{s}^{\prime}\) its dual equipped with the topology \(\sigma(F^{\prime},F)\) . We denote by \(L_{F_{s}^{\prime}}^{\infty}\) the vector space of mappings f of T into \(F_{s}^{\prime}\) , such that f is scalarly \(\mu\) -measurable and is equal scalarly locally almost everywhere (for \(\mu\) ) to a mapping of T into an equicontinuous subset of \(F^{\prime}\) . We denote by \(L_{F_{s}^{\prime}}^{\infty}\) the quotient space of \(L_{F_{s}^{\prime}}^{\infty}\) by the space of scalarly locally \(\mu\) -negligible mappings of T into \(F_{s}^{\prime}\) .

When F satisfies the hypotheses of §1, No.~5, Prop. 13, the functions in \(\mathcal{L}_{\mathrm{F}_s^\prime}^\infty\) are \(\mu\) -measurable for the weak topology \(\sigma(\mathrm{F}',\mathrm{F})\) , but are not necessarily measurable for the strong topology on \(\mathrm{F}'\) , even if F is a Banach space (§1, Exer. 17). Under the same conditions, the scalarly locally \(\mu\) -negligible mappings of T into \(\mathrm{F}_s^\prime\) are identical to the locally \(\mu\) -negligible mappings of T into \(\mathrm{F}_s^\prime\) (§1, No.~1, Remark 2).

When F is a separable normed space, the elements of \(L_{F_{s}^{\prime}}^{\infty}\) are the mappings f of T into \(F_{s}^{\prime}\) such that f is scalarly \(\mu\) -measurable and \(|f|\) is bounded in measure; one can then define a normed space structure on the space \(L_{F_{s}^{\prime}}^{\infty}\) , by equipping it with the norm \(N_{\infty}\) (Ch. IV, §6, No.~3).

Lemma 3. --- Let F be a Hausdorff locally convex space, f an element of \(\mathcal{L}_{\mathrm{F}_s^\prime}^\infty\) . For every \(\mathbf{z} \in \mathrm{F}\) , one has \(\langle \mathbf{z}, \mathbf{f} \rangle \in \mathcal{L}^\infty\) , and the linear mapping \(z \mapsto \pi(\langle \mathbf{z}, \mathbf{f} \rangle)\) of F into \(\mathrm{L}^\infty\) is continuous; if, moreover, F is a normed space, then \(\mathrm{N}_\infty(\langle \mathbf{z}, \mathbf{f} \rangle) \leqslant |\mathbf{z}| \cdot \sup_{t \in \mathrm{T}} |\mathbf{f}(t)|\) .

It is clear by definition that \(\langle z,f\rangle\) is \(\mu\) -measurable and bounded in measure; replacing if necessary f by a function belonging to the same class of \(L_{F_{s}}^{\infty}\) , we can suppose in addition that \(\mathbf{f}(T)\subset V^{\circ}\) , where V is a balanced convex neighborhood of 0 in F (which does not modify \(\langle z,f\rangle\) except on a locally negligible set, depending on z). Then the relation \(z\in V\) implies that \(|\langle z,f(t)\rangle|\leqslant1\) for all \(t\in T\) , which proves the continuity of \(\mathbf{z}\mapsto\pi(\langle\mathbf{z},\mathbf{f}\rangle)\) . The final assertion is obvious.

Lemma 4. --- Let F be a Hausdorff locally convex space, f an element of \(\mathcal{L}_{\mathrm{F}_s}^\infty\) . For every numerical function \(g \in \overline{\mathcal{L}}^1\) , the function gf is scalarly essentially \(\mu\) -integrable and \(\int g\mathbf{f} d\mu \in \mathrm{F}'\) .

For every \(z \in F\) , \(\langle z, f \rangle\) belongs to \(L^{\infty}\) , whence the first assertion. One can suppose moreover, without modifying \(\int g f d\mu\) , that \(\mathbf{f}(T) \subset V^{\circ}\) , where V is a balanced convex neighborhood of 0 in F. Then the relation \(z \in V\) implies \(|\langle z, f(t) \rangle| \leqslant 1\) for all \(t \in T\) , whence \(|\langle z, \int g f d\mu \rangle| = |\int \langle z, f \rangle g d\mu| \leqslant \overline{N}_{1}(g)\) , which proves that \(\int g f d\mu \in F'\) .

THEOREM 1. --- Let F be a Hausdorff locally convex space that contains a countable dense subset. For every function \(\mathbf{f} \in \mathcal{L}_{\mathrm{F}_s}^\infty\) and every \(\mathbf{z} \in \mathbf{F}\) , let \(v_{\mathbf{f}}(\mathbf{z}) = \pi (\langle \mathbf{z}, \mathbf{f} \rangle) \in \mathrm{L}^{\infty}\) ; the mapping \(\mathbf{f} \mapsto v_{\mathbf{f}}\) defines, by passage to the quotient, a linear bijection of \(L_{F^{\prime}s}^{\infty}\) onto the space \(\mathcal{L}(F;L^{\infty})\) of continuous linear mappings of F into \(L^{\infty}\) . If F is a normed space, this bijection is an isometry.

In view of Lemma 3, the first assertion will be demonstrated if one proves that for every continuous mapping \(u\) of \(F\) into \(L^{\infty}\) , there exists a function \(\mathbf{f} \in \mathcal{L}_{F_s'}^{\infty}\) such that \(\pi(\langle \mathbf{z}, \mathbf{f} \rangle) = u(\mathbf{z})\) for all \(\mathbf{z} \in F\) , and that the class of \(\mathbf{f}\) in \(L_{F_s'}^{\infty}\) is uniquely determined by this condition. The second point is immediate, because if \(\pi(\langle \mathbf{z}, \mathbf{f} \rangle) = \pi(\langle \mathbf{z}, \mathbf{f}_1 \rangle)\) for all \(\mathbf{z} \in F\) , then \(\mathbf{f}_1 - \mathbf{f}\) is scalarly locally negligible. On the other hand, there exists a lifting \(\rho\) of \(u(F)\) into \(\mathcal{L}^{\infty}\) (Lemma 2). For every \(t \in T\) , the mapping \(\mathbf{z} \mapsto \rho(u(\mathbf{z}))(t)\) is a continuous linear form on \(F\) , that is, an element \(f(t)\) of \(F'\) . The function \(\mathbf{f}\) is scalarly \(\mu\) -measurable since \(\langle \mathbf{z}, \mathbf{f} \rangle = \rho(u(\mathbf{z})) \in \mathcal{L}^{\infty}\) for every \(\mathbf{z} \in F\) ; one has \(\pi(\langle \mathbf{z}, \mathbf{f} \rangle) = u(\mathbf{z})\) ; finally, for every \(t \in T\) and every \(\mathbf{z}\) belonging to the inverse image \(V\) under \(u\) of the unit ball of \(L^{\infty}\) ,

\[
\left| \langle \mathbf {z}, \mathbf {f} (t) \rangle \right| = \left| \rho (u (\mathbf {z})) (t) \right| \leqslant N _ {\infty} (u (\mathbf {z})) \leqslant 1,
\]

which shows that \(\mathbf{f}(t) \in \mathrm{V}^{\circ}\) for all \(t \in \mathrm{T}\) .

If, moreover, F is a normed space, the foregoing shows that

\[
\sup _ {t \in \mathrm{T}} | \mathbf {f} (t) | \leqslant \| u \|.
\]

But on the other hand (Lemma 3), \(\mathrm{N}_{\infty}(u(\mathbf{z})) \leqslant |\mathbf{z}| \cdot \sup_{t \in \mathrm{T}} |\mathbf{f}(t)|\) , and this inequality continues to hold when f is modified arbitrarily on a locally negligible set. It follows that \(\|u\| = \mathrm{N}_{\infty}(|\mathbf{f}|)\) .

COROLLARY 1. --- Let F be a Hausdorff locally convex space containing a countable dense subset. For every function \(\mathbf{f} \in \mathcal{L}_{\mathrm{F}_s}^\infty\) , every \(\mathbf{z} \in \mathrm{F}\) and every function \(g \in \mathcal{L}^1\) , let \(\Phi_{\mathbf{f}}(\mathbf{z}, \widetilde{g}) = \int \langle \mathbf{z}, \mathbf{f}(t) \rangle g(t) d\mu(t)\) . The mapping \(\mathbf{f} \mapsto \Phi_{\mathbf{f}}\) defines, by passage to the quotient, a linear bijection of \(\mathrm{L}_{\mathrm{F}_s}^\infty\) onto the space \(\mathcal{B}(\mathrm{F}, \mathrm{L}^1)\) of continuous bilinear forms on \(\mathrm{F} \times \mathrm{L}^1\) . If \(\mathrm{F}\) is a normed space, this bijection is an isometry.

One can suppose that \(\mathbf{f}(\mathrm{T})\) is an equicontinuous subset of \(F'\) . It is then clear that \(\Phi_{f}\) is separately continuous, and, with the notations of Th. 1 and the Appendix, one has (taking into account the fact that \(L^{\infty}\) is the dual of \(L^{1}\) (Ch. V, §5, No.~8, Th. 4)) \(^{l}\Phi_{f}=v_{f}\) . The corollary then follows from Th. 1 above and from the Appendix, No.~1, Prop. 1 and its corollary.

COROLLARY 2 (Dunford--Pettis theorem). --- Let F be a Hausdorff locally convex space containing a countable dense subset. For every function \(f \in L_{F_{s}^{\infty}}\) and every function \(g \in L^{1}\) , let \(w_{\mathbf{f}}(\widetilde{g}) = \int g\mathbf{f} \, d\mu \in \mathbf{F}'\) (Lemma 4). The mapping \(f \mapsto w_{f}\) defines, by passage to the quotient, a linear bijection of \(L_{F_{s}^{\prime}}^{\infty}\) onto the space \(\mathcal{R}(L^{1}, F^{\prime})\) of linear mappings u of \(L^{1}\) into \(F^{\prime}\) such that the image under u of the unit ball of \(L^{1}\) is an equicontinuous subset of \(F^{\prime}\) . If F is a normed space (in which case \(\mathcal{R}(L^{1}, F^{\prime})\) is the space of continuous linear mappings of \(L^{1}\) into the strong dual of F), the bijection \(f \mapsto w_{f}\) is an isometry.

Taking into account the fact that \(L^{\infty}\) is the dual of \(L^{1}\) , this follows from the preceding corollary and from the Appendix, No.~1, Prop. 1 and its corollary.

Remark. --- It is clear that the mappings \(u \in \mathcal{R}(L^{1}, F')\) are continuous for every S-topology on \(F'\) (S a covering of F by bounded subsets). Conversely, if F is assumed moreover to be barreled, then every continuous linear mapping of \(L^{1}\) into \(F'\) equipped with an S-topology transforms the unit ball of \(L^{1}\) into a bounded subset of \(F'\) , which is therefore equicontinuous (TVS, III, §4, No.~2, Th. 1).

COROLLARY 3. --- Let F be a Hausdorff locally convex space containing a countable dense subset, m a vectorial measure on T with values in the weak dual F' of F. If the image under m of the set B of functions g in \(\mathcal{K}(T)\) such that \(\mu(|g|) \leqslant 1\) is contained in a closed and convex equicontinuous subset H' of F', then m has base \(\mu\) and there exists a density f of m with respect to \(\mu\) such that f(t) ∈ H' for all t ∈ T.

The hypothesis implies that m is continuous when \(\mathcal{K}(T)\) is equipped with the topology induced by that of \(\mathcal{L}^{1}(\mu)\) (defined by the semi-norm \(N_{1}\) ); it may therefore be extended to a continuous linear mapping u of \(\mathcal{L}^{1}(\mu)\) into the completion G of the weak dual of F; but since \(H'\) is a compact subset of G and the image under u of the set \(\overline{B}\) of \(f \in L^{1}\) such that \(\mathrm{N}_{1}(f) \leqslant 1\) is contained in the closure of \(H'\) in G, one has \(u(\overline{\mathrm{B}}) \subset \mathrm{H}'\) , therefore u maps \(L^{1}\) into \(F'\) . Since the relation \(\mathrm{N}_{1}(f) \leqslant \varepsilon\) implies that \(u(f) \in \varepsilon\mathrm{H}'\) , one has \(u(g) = 0\) if g is \(\mu\) -negligible, and Cor. 2 may therefore be applied to the mapping of \(L^{1}\) into \(F'\) obtained from u by passing to the quotient; whence the corollary.

COROLLARY 4. --- Let F be a separable normed space, and m a measure on T with values in the strong dual F', majorizable for the norm of F'. Then m is a measure with base \textbar m\textbar, and if \(m = f \cdot |m|\) then \(|f(t)| = 1\) locally almost everywhere for \textbar m\textbar.

By hypothesis, for every \(\mathbf{z} \in \mathbf{F}\) such that \(|\mathbf{z}| \leqslant 1\) , one has \(|\langle \mathbf{z}, \mathbf{m}(g) \rangle| \leqslant |\mathbf{m}|(|g|)\) for every function \(g \in \mathcal{K}(\mathrm{T})\) , consequently \(|\mathbf{m}(g)| \leqslant |\mathbf{m}|(|g|)\) (TVS, IV, §2, No.~4, formula (1)). Since every ball in \(\mathbf{F}'\) is equicontinuous, Cor. 3 is applicable and shows that \(\mathbf{m}\) has base \(|\mathbf{m}|\) ; moreover, if \(\mathbf{m} = \mathbf{f} \cdot |\mathbf{m}|\) then \(\mathbf{f}\) is \(|\mathbf{m}|\) -measurable for \(\sigma(\mathrm{F}', \mathrm{F})\) (§1, No.~5, Prop. 13) and \(|\mathbf{m}| = |\mathbf{f}| \cdot |\mathbf{m}|\) (No.~4, Prop. 8), which proves the corollary (Ch. V, §5, No.~3, Cor. 2 of Prop. 3).

If this corollary is applied to the case that \(\mathbf{F}\) if finite-dimensional, one recovers as a special case the first part of Prop. 9.

